% !PRINTED 610
\weberpage{610}
%<<CONT
equation to admit solution by radicals. We therefore derive, at first under the
assumption that the degree $n$ of the equation is a prime, another criterion first set
up by \gesperrt{Galois}.

Now let $f(x) = 0$ be an irreducible equation of degree $n$, and $n$ a prime greater
than 2. If $f(x)$ is to be algebraically solvable, then by \S.\,176 its group $P$ must
be metacyclic, i.e. there must be a chain of groups
\weq{1}{P,\ P_1,\ P_2 \ldots P_{\mu-1},\ 1}
each of which is a normal divisor of the one immediately preceding, with prime index.
By successive adjunction of roots of cyclic equations of prime degree the group of the
equation is reduced from $P$ to $P_1, P_2 \ldots P_{\mu-1}, 1$.

The function $f(x)$ itself, as was already shown in \S.\,176, cannot be reduced before
the last adjunction has been made, whereupon it splits into $n$ linear factors. The
degree of the last cyclic equation, and hence the degree of the next-to-last group
$P_{\mu-1}$, must therefore by \S.\,158, 3. be equal to $n$.

The degree of an element of a group is always a divisor of the degree of the group, and
hence $P_{\mu-1}$ can contain, besides the identical one, only permutations of order
$n$. It is therefore identical with the period $1, \pi, \pi^2 \ldots \pi^{n-1}$ of one
of its elements. From this it follows by \S.\,153 that $\pi$ must be a cyclic
permutation of all $n$ figures, and we can choose the designation of the roots of
$f(x)$ so that
\[ \pi = (0,\ 1,\ 2 \ldots n-1) \]
holds. If we denote the roots by $x_z$, and stipulate that $x_z = x_{z'}$ shall hold
when $z \equiv z' \, (\mathrm{mod}\ n)$, then by $\pi$ every $z$ passes into $z+1$, by
$\pi^2$ into $z+2$, and so on, and we can therefore say that the group $P_{\mu-1}$
consists of the substitutions for $z$
\weq{2}{(z,\ z+b), \qquad b = 0,\ 1,\ 2 \ldots n-1}
Every irreducible metacyclic equation of prime degree can therefore be transformed by
adjunction of radicals into a cyclic equation.

The substitutions $(z, z+b)$ form a special case of the more general linear
substitution
\weq{3}{(z,\ az+b),}
in which $a$ and $b$ are fixed numbers, which may also be reduced according to the
modulus $n$, but where for $a$ the value 0 of course
%CONT>>
